% ================== Chapitre 6 : controles de recette =======================
\chapter{Contrôles de recette de la tâche}
\label{ch:recette}

\section{Critères d'acceptation et commandes associées}

Les contrôles ci-dessous définissent ce que signifie exactement, pour T005, que
le service DNS est en place. Ils portent sur le service lui-même et non sur le
comportement d'ensemble du datacenter. Conformément à la règle du corpus, leur
statut initial est \statut{non exécuté} et aucun résultat n'est présenté comme
obtenu.

\begin{longtable}{@{}C{0.09}R{0.19}R{0.37}R{0.35}@{}}
\caption{Contrôles de recette de T005 ; statut initial de l'ensemble :
  \statut{non exécuté}}\label{tab:controles}\\
\hautfilet
\textbf{Code} & \textbf{Objet} & \textbf{Précondition et action} &
\textbf{Résultat attendu}\\
\medfilet
\endfirsthead
\hautfilet
\textbf{Code} & \textbf{Objet} & \textbf{Précondition et action} &
\textbf{Résultat attendu}\\
\medfilet
\endhead
\basfilet
\endfoot
C01 & Réponse autoritative directe &
  \texttt{dig +norecurse @10.30.4.10 ns1.gandal.internal A} &
  \opt{NOERROR}, indicateur \opt{aa} présent, réponse \ip{10.30.4.10}\\
C02 & Cohérence entre les deux autoritatifs &
  \texttt{dig +short @10.30.4.10 tenant-007.gandal.internal SOA} puis la même
  requête sur \ip{10.30.4.11} &
  Numéros de série identiques après réception du \opt{NOTIFY}\\
C03 & Résolution inverse d'une machine tenant &
  \texttt{dig +short @10.30.4.12 -x 10.30.70.37} &
  \zone{vm-0042.tenant-007.\allowbreak gandal.internal.}\\
C04 & Transfert de zone refusé à un tiers &
  \texttt{dig @10.30.4.10 gandal.internal AXFR} depuis une source autre que ns2 &
  Transfert refusé ; le transfert signé demandé par ns2 reste possible\\
C05 & Isolation entre tenants &
  depuis \ip{10.30.70.37} : \texttt{dig @10.30.4.12
  gw.tenant-008.\allowbreak gandal.internal A} &
  \opt{REFUSED}, sans divulgation de l'existence du nom\\
C06 & Isolation de l'infrastructure &
  depuis \ip{10.30.70.37} : \texttt{dig @10.30.4.12 xo.infra.gandal.internal A} &
  \opt{REFUSED} ; la même requête depuis le pool VPN administrateur aboutit\\
C07 & Récursion autorisée et bornée &
  depuis \ip{10.30.70.37} : \texttt{dig @10.30.4.12 example.org A} ; puis la
  même requête depuis une source hors \opt{allow-from} &
  \opt{NOERROR} dans le premier cas, absence de service dans le second\\
C08 & Autoritatif inatteignable depuis un tenant &
  depuis \ip{10.30.70.37} : \texttt{dig +time=2 +tries=1 @10.30.4.10
  gw.tenant-007.\allowbreak gandal.internal A} &
  Absence de réponse, refus journalisé sur la passerelle\\
C09 & Bascule sur perte d'une instance &
  Arrêt de ns1 puis requête sur ns2 ; arrêt de rec1 puis requête depuis un
  client dont le \opt{resolv.conf} liste \ip{.12} et \ip{.13} &
  Résolution maintenue dans les deux cas ; publication suspendue pendant
  l'arrêt de ns1\\
C10 & Idempotence de la publication &
  Émission deux fois du même appel du listing~\ref{lst:publication} &
  État final identique, aucun doublon, aucun second transfert si la lecture
  préalable est implémentée\\
C11 & Mode dégradé de publication &
  Arrêt de l'API autoritative pendant une publication &
  Opération maintenue en état observable, adresse non rendue au pool, alerte
  levée\\
C12 & Amorçage sans DNS &
  Depuis H1, accès à XO et au bastion avec le service DNS arrêté &
  Accès obtenu par le socle de reprise du listing~\ref{lst:hosts}\\
C13 & Comportement sous contrainte de MTU &
  \texttt{dig +dnssec +bufsize=1232} sur une réponse volumineuse, puis
  \texttt{dig +tcp} &
  Troncature suivie d'une reprise en TCP, sans perte silencieuse\\
C14 & Journaux et métriques en lecture pull &
  Lecture du point de métriques depuis \ip{10.30.5.10}, puis tentative
  d'initiation depuis une instance DNS vers Monitoring &
  Lecture réussie, initiation inverse refusée, horodatage en temps universel\\
C15 & Discrétion du service &
  \texttt{dig @10.30.4.10 version.bind CH TXT} & Réponse anonymisée\\
\end{longtable}

\section{Transmission à la campagne de validation}

Plusieurs scénarios du corpus dépassent le périmètre de cette tâche parce qu'ils
engagent d'autres lots ou mesurent des cibles globales. Ils sont rappelés ici
avec la contribution que T005 leur apporte, afin que la campagne de validation
puisse les reprendre sans redéfinir le contexte.

\begin{table}[H]
\centering
\small
\caption{Scénarios transmis à la campagne de validation}\label{tab:transmis}
\begin{tabular}{@{}R{0.16}R{0.34}R{0.50}@{}}
\hautfilet
\textbf{Scénario} & \textbf{Objet} & \textbf{Contribution de T005}\\
\medfilet
S04, document 03 &
  Requêtes DNS vers l'infrastructure et vers un tenant voisin depuis le VNet
  USER &
  C05 et C06 en constituent la vérification unitaire ; S04 y ajoute le parcours
  complet depuis une machine réellement provisionnée.\\
S08, document 03 &
  Restauration depuis une sauvegarde indépendante &
  La procédure et le contenu de la sauvegarde sont fournis à la
  section~\ref{sec:sauvegarde} ; la mesure du délai reste à produire.\\
G04, document 01 &
  Coupure de XO puis d'un service DNS &
  Le comportement attendu est décrit à la section~\ref{sec:suppression} :
  continuité du trafic existant, suspension des opérations nouvelles.\\
V09, document 02 &
  Création de cinquante VNets et cinq cents baux &
  La génération des cent treize zones est automatisée par le
  listing~\ref{lst:zones} ; la tenue en charge du service reste à mesurer.\\
NET-NFR-003 &
  Reconstruction en moins de soixante minutes &
  Éléments de reprise fournis aux sections~\ref{sec:amorcage}
  et~\ref{sec:sauvegarde} ; la cible est à évaluer et non à présumer.\\
\basfilet
\end{tabular}
\end{table}
